\section{Introduction}

The theory of Frobenius manifolds, having its origin in theoretical physics, has deep interrelations with apparently very dif\/ferent areas of mathematics: Gromov--Witten invariants and quantum cohomology, integrable systems, singularity, deformation of f\/lat connections etc.
We discuss a~new aspect of this fruitful and fast developing theory: its relations with the classical chapter of dif\/ferential geometry, namely the web theory.

The notion of Frobenius manifold, introduced by B.~Dubrovin (see~\cite{Df} and~\cite{Mf}), is a geometric translation of the theory of WDVV-equations that arise originally in the physical context of two-dimensional topological f\/ield theory.
\begin{definition}\label{defFrob}
A Frobenius manifold is a complex analytic manifold $M$ equipped with the following
analytic objects:
\begin{enumerate}\itemsep=0pt
\item[1)] a commutative and associative multiplication on $T_pM$,
\item[2)] an invariant non-degenerate f\/lat inner product: $\langle u \cdot v, w \rangle = \langle u, v \cdot w \rangle $,
\item[3)] a constant unity vector f\/ield $e$: $\nabla e=0$, $e \cdot v=v$ $\forall\, v\in TM$,
\item[4)] a linear Euler vector f\/ield $E$: $\nabla (\nabla E)=0$,
\end{enumerate}
satisfying the following conditions:
\begin{itemize}\itemsep=0pt
\item the f\/low of $E$ re-scales the
multiplication and the inner product,
\item 4-tensor $(\nabla_z c)(u,v,w)$ is symmetric in
$u$, $v$, $w$, $z$, where
\[
c(u,v,w):=\langle u\cdot v,w\rangle.
\]
\end{itemize}
In this def\/inition, the symbol $\nabla$ stands for the Levi-Civita connection of the inner product $\langle\ ,\ \rangle$.
\end{definition}
The following geometric construction of $3$-web via Frobenius $3$-fold was proposed in~\cite{Aw}.
Consider a {\it massive} Frobenius $3$-fold $M$, which means that the algebra $T_pM$ is semi-simple for each $p\in U$ for some open set $U\subset M$. Then $T_pM$ is the direct
product of one-dimensional algebras spanned by idempotents
\begin{equation}\label{mult}
T_pM=\mathbb C \{e_1\}\otimes \mathbb C \{e_2\}\otimes \mathbb C\{e_3\}, \qquad e_i\cdot e_j=\delta_{ij}e_i.
\end{equation}
In this setting, the unity vector f\/ield is
\[
e=e_1+e_2+e_3.
\]
Let $S$ be a surface transverse to the unit vector f\/ield $e$, then 3
planes spanned by $\{e,e_i\}$ cut 3 directions on $T_pS$ (see Fig.~\ref{booklet} on the left).
The integral curves of these direction f\/ields build a {\it flat} (or {\it hexagonal}) {\it $3$-web}, i.e., at the points with pairwise distinct web directions, this web is locally biholomorphic to 3 families of parallel lines in the plane.
Due to the existence of {\it canonical Dubrovin} coordinates (see~\cite{Di}), the distributions $\{e,e_i\}$ are integrable, integral surfaces of each of the distributions being formed by the integral curves of the unity vector f\/ield~$e$. Thus for each point in $M$ there are 3 integral surfaces intersecting along such a curve. These surfaces cut $S$ along the constructed web. This justif\/ies the following def\/inition (see also Fig.~\ref{booklet} on the right).
\begin{figure}[th]\centering
\includegraphics[width=40mm]{Agafonov-Fig1a} \qquad\quad
\includegraphics[width=40mm]{Agafonov-Fig1b}
%\includegraphics[width=80mm]{aaa}
%\hspace{-0.5cm} \epsfig{file=idgif_eps.eps,width=80mm}
%\qquad \put(1,8){ \epsfig{file=booklet.eps,width=40mm}}
\caption{Construction of a booklet $3$-webs from Frobenius $3$-folds.}
%$e_i$ \quad $T_p S$ \quad $S$ \quad $e$
\label{booklet}
\end{figure}

\begin{definition}
The constructed web is called a booklet $3$-web.
\end{definition}
The web directions of a booklet $3$-web are well-def\/ined everywhere (see~\cite{Aw}), maybe with multiplicity.
We call a point {\it regular or non-singular} if the web directions at this point are pairwise distinct. The locus of singular points is called the discriminant curve.
By def\/inition, a hexagonal $3$-web does not have any local invariants at regular points. Its ``personality'' is encoded in the behavior at \emph{singular points}, where at least 2 web directions coincide.
